\section{Оценка}
\label{sec:evaluation}

Главите~\ref{sec:enforcement} и~\ref{sec:case-study} описаха как гаранцията се
постига и как беше затворен един пропуск в нея. Настоящата глава проверява
гаранцията по метода от точка~\ref{sec:tm-method} и съпоставя резултатите с
критериите от таблица~\ref{tab:tm-criteria}. Наблюденията се представят първо, а
тълкуването им следва отделно, за да може читателят да прецени едното независимо
от другото.

\subsection{Среда и възпроизводимост}
\label{sec:eval-env}

Всички измервания са извършени с изходния код във версия \code{a17f4e7}, при
следната среда:

\begin{itemize}
    \item операционна система Ubuntu 24.04.4 LTS под WSL2, ядро
          \code{6.6.114.1-microsoft-standard-WSL2};
    \item компилатори GCC 14.2.0 и Clang 18.1.3, стандарт C++23;
    \item SQLite 3.45.1, база данни в паметта, нова за всяко изпълнение;
    \item CMake 3.28.3.
\end{itemize}

Наборът се изгражда и изпълнява с опцията \code{ORM\_ENABLE\_SECURITY}, както е
описано в приложение~А. Всеки резултат се записва в каталога
\code{security/evidence/} в изходния код, а не в каталога за изграждане. Така
всяко изпълнение е преглеждаема разлика в системата за контрол на версиите и
всяко число в тази глава може да бъде проследено до файл. От версия
\code{a17f4e7} нататък наборът се изпълнява и при всяка промяна в системата за
непрекъсната интеграция, като изпълнението за самата версия \code{a17f4e7} е
успешно.

\subsection{Корпус от полезни товари}
\label{sec:eval-corpus}

Корпусът съдържа 66 полезни товара от 11 публични източника. Всеки ред съдържа
препратка към източника си и нито един полезен товар не е съставен за целите на
работата. Таблица~\ref{tab:eval-corpus} показва разпределението.

\begin{table}[H]
\centering
\caption{Състав на корпуса}
\label{tab:eval-corpus}
\begin{tabularx}{\textwidth}{L{3.0cm}L{2.0cm}Y}
\toprule
\textbf{Клас} & \textbf{Брой} & \textbf{Диалекти} \\
\midrule
SQL & 36 & общи 16, MySQL 8, PostgreSQL 6, SQLite 6 \\
NoSQL (MongoDB) & 15 & инжекция на оператори \\
Cypher (Neo4j) & 6 & \\
CQL (Cassandra) & 4 & \\
Redis & 5 & инжекция на протоколно ниво \\
\midrule
Общо & 66 & \\
\bottomrule
\end{tabularx}
\end{table}

Разпределението е неравномерно и това е следствие от правилото полезен товар да
се включва само ако е открит дословно в източник. Три полезни товара са
публикувани от източника с незапълнено поле, например \code{[RANDNUM]}. Те се
изпълняват точно както са публикувани, защото запълването на полето би означавало
съставяне на нов текст. Корпусът изрично изброява и пропуските си, сред които
липсата на проверен полезен товар с \code{UNION} и с натрупани заявки за SQLite,
както и с \code{UNION} за Cypher.

\subsection{Отрицателни компилационни тестове}
\label{sec:eval-compile}

Наборът съдържа 18 компилационни теста: 15 програми, които трябва да бъдат
отхвърлени, и 3 контролни случая, които трябва да се компилират.
Таблица~\ref{tab:eval-compile} показва резултата при двата компилатора.

\begin{table}[H]
\centering
\caption{Резултати от компилационните тестове}
\label{tab:eval-compile}
\begin{tabularx}{\textwidth}{L{6.3cm}L{1.6cm}L{1.9cm}YY}
\toprule
\textbf{Случай} & \textbf{Ниво} & \textbf{Очакване} & \textbf{GCC} & \textbf{Clang} \\
\midrule
00 документирана заявка        & контрола & компилира & да & да \\
14 сравнение с \code{nullptr}  & контрола & компилира & да & да \\
15 сравнение на две колони      & контрола & компилира & да & да \\
\midrule
06 низ към целочислен заместител & 1 & отхвърля & да & да \\
01 име на колона по време на изпълнение & 2 & отхвърля & да & да \\
16 име на колона в \code{std::string} & 2 & отхвърля & да & да \\
02 име на таблица по време на изпълнение & 2 & отхвърля & да & да \\
04 колона, избрана по време на изпълнение & 2 & отхвърля & да & да \\
17 две колони в една променлива & 2 & отхвърля & да & да \\
03 текст на заявка към \code{execute} & 3 & отхвърля & да & да \\
05 условие като текст & 3 & отхвърля & да & да \\
\midrule
07 литерал, \code{execute} & литерал & отхвърля & да & да \\
08 литерал, \code{operator<<} & литерал & отхвърля & да & да \\
09 литерал, \code{find\_one} & литерал & отхвърля & да & да \\
10 литерал, \code{prepare} & литерал & отхвърля & да & да \\
11 литерал в \code{SET} & литерал & отхвърля & да & да \\
12 литерал в \code{DELETE} & литерал & отхвърля & да & да \\
13 литерал, константа & литерал & отхвърля & да & да \\
\bottomrule
\end{tabularx}
\end{table}

„Да“ в таблицата означава, че програмата е дала очаквания резултат и, при
отхвърляне, че диагностиката съответства на образеца за проверяваната причина.
Диагностиките са запазени поотделно за всеки компилатор в
\code{security/evidence/compile/}.

Два от тестовете бяха уточнени след първото изпълнение и това има значение за
тълкуването. Първоначалните варианти на случаи 01 и 04 бяха отхвърлени, но по
причина, различна от проверяваната, и затова се отчетоха като неуспешни. Случай
01 опитваше да използва \code{std::string} като име на колона и беше отхвърлен
поради несъответствие на типа, преди въпросът за константността изобщо да
възникне. Той беше разделен на случай 01, който подава стойност с правилния тип,
но известна едва по време на изпълнение, и случай 16, който запазва варианта с
\code{std::string}. Случай 04 опитваше да избере между две колони и беше
отхвърлен със съобщението, че двата указателя към член имат различни типове.
Това разкри по-силно свойство от проверяваното: понеже името на колоната е част
от типа на свойството, две колони с един и същ тип на стойността (тук и двете
\code{int}) са различни типове, и една променлива не може да обхваща и двете.
То е отделено в случай 17, а случай 04 проверява само константността.
Записът от първото изпълнение е запазен в
\code{security/evidence/compile\_before\_fix.txt}.

\subsection{Диференциален тест със SQLite}
\label{sec:eval-differential}

От корпуса към SQLite са приложими 22 полезни товара: общите 16 и специфичните
за SQLite 6. Таблица~\ref{tab:eval-diff-summary} обобщава резултата.

\begin{table}[H]
\centering
\caption{Диференциален тест със SQLite, 22 полезни товара}
\label{tab:eval-diff-summary}
\begin{tabularx}{\textwidth}{L{5.6cm}YY}
\toprule
\textbf{Резултат} & \textbf{Наивна заявка} & \textbf{Библиотеката} \\
\midrule
Компрометирана: върнати редове & 4 & 0 \\
Компрометирана: промяна на състоянието & 0 & 0 \\
Компрометирана: принудително изпълнение & 0 & 0 \\
Синтактично нарушена & 10 & 0 \\
Без наблюдаем ефект & 8 & 22 \\
\bottomrule
\end{tabularx}
\end{table}

Таблица~\ref{tab:eval-diff-technique} показва резултата на наивната заявка по
техники.

\begin{table}[H]
\centering
\caption{Резултат на наивната заявка по техники}
\label{tab:eval-diff-technique}
\begin{tabularx}{\textwidth}{L{4.2cm}L{1.6cm}Y}
\toprule
\textbf{Техника} & \textbf{Брой} & \textbf{Резултат} \\
\midrule
тавтология & 2 & и двата компрометират (3 и 1 върнати реда) \\
\code{UNION} & 2 & и двата компрометират (по 1 върнат ред) \\
сляпа атака по време & 6 & и шестте нарушават синтаксиса \\
натрупана заявка & 1 & нарушава синтаксиса \\
сляпа атака по истинност & 3 & 2 без ефект, 1 нарушава синтаксиса \\
коментар & 3 & 2 без ефект, 1 нарушава синтаксиса \\
изброяване & 2 & 1 без ефект, 1 нарушава синтаксиса \\
предизвикване на грешка & 1 & без ефект \\
заобикаляне на екраниране & 2 & без ефект \\
\bottomrule
\end{tabularx}
\end{table}

Наблюденията са следните. Наивната заявка е компрометирана от 4 от 22 полезни
товара. И четирите са върнали редове, които заявката не би трябвало да върне.
Единият от двата полезни товара с \code{UNION}, \code{SQL-003}, е върнал реда с
потребителското име и паролата от таблицата \code{users}, тоест е извлякъл
съхранена идентификационна информация. Другият, \code{SQL-004}, е върнал ред с
избрани от нападателя стойности, техника за заобикаляне на удостоверяване. През
библиотеката и 22-та полезни товара не са имали наблюдаем ефект. Пълният резултат
по полезни товари е в \code{security/evidence/differential\_sqlite.csv}.

\subsection{Тълкуване}
\label{sec:eval-interpretation}

Резултатът на библиотеката, 22 от 22 без ефект, придобива смисъл едва на фона на
резултата на наивната заявка. Четирите компрометирания показват, че корпусът и
постановката са способни да разкрият уязвимост, тоест че нулата при библиотеката
не се дължи на безсилен корпус.

Десетте синтактично нарушени заявки заслужават отделен коментар. Те не са
неуспех на атаката в смисъла на защитата. Полезният товар е достигнал граматиката
на SQLite и е нарушил заявката, но не е постигнал ефект. Причината е проверена за
всеки случай. Пет от шестте слепи атаки по време използват конструкции на други
системи, които SQLite не поддържа: \code{SLEEP} на MySQL, \code{WAITFOR DELAY} на
SQL Server и \code{dbms\_pipe} на Oracle. Шестата, \code{SQL-034}, е специфична за
SQLite, но е публикувана с незапълнени полета и затова не е синтактично валидна.
Натрупаната заявка \code{SQL-017} извиква процедурата \code{xp\_cmdshell} на SQL
Server. Срещу библиотеката същите полезни товари не достигат граматиката изобщо:
те пристигат като стойност на свързан параметър и се сравняват като низ.

Осемте полезни товара без ефект при наивната заявка са предимно слепи атаки и
опити за заобикаляне на екраниране. Слепата атака по дефиниция не дава видим
резултат при едно изпълнение: тя разчита на разлика между две изпълнения. Затова
оракулът, който вижда само ефекти, не я отчита. Както е посочено в
точка~\ref{sec:tm-not-proven}, броят на компрометиранията на наивната заявка е
долна граница.

\subsection{Сървърните хранилища}
\label{sec:eval-servers}

Диференциалният тест е изпълнен със SQLite. За останалите шест хранилища
приложимите полезни товари са: PostgreSQL 22, MySQL 24, MongoDB 15, Neo4j 6,
Cassandra 4 и Redis 5. Изпълнението им изисква работещи сървъри, които се
стартират в контейнери. В средата, в която са извършени измерванията за
настоящата версия на работата, услугата за контейнери не беше налична и тези
изпълнения не са извършени. За тези хранилища гаранцията е потвърдена чрез
анализа на изходния код в точка~\ref{sec:enf-backends}, но не и чрез измерване.

Конекторът за Neo4j има допълнително ограничение. Образът за контейнери, използван
в системата за непрекъсната интеграция, не съдържа клиентската библиотека
\code{libneo4j-client} и не стартира сървър Neo4j. Затова за Neo4j и
съществуващите интеграционни тестове на библиотеката не се изпълняват. Измерването
му изисква разширяване на образа и е посочено като бъдеща работа.

\subsection{Съпоставка с критериите за успех}
\label{sec:eval-criteria}

Таблица~\ref{tab:eval-criteria} съпоставя резултатите с критериите от
таблица~\ref{tab:tm-criteria}.

\begin{table}[H]
\centering
\caption{Съпоставка с критериите за успех}
\label{tab:eval-criteria}
\begin{tabularx}{\textwidth}{L{1.0cm}L{2.8cm}Y}
\toprule
\textbf{№} & \textbf{Оценка} & \textbf{Основание} \\
\midrule
К1 & изпълнен за SQLite & 22 от 22 полезни товара без ефект през библиотеката; за сървърните хранилища само чрез анализ \\
К2 & изпълнен & 15 от 15 отхвърляния по проверяваната причина и 3 от 3 контролни случая, при GCC и при Clang \\
К3 & изпълнен за SQLite & наивната заявка е компрометирана от 4 полезни товара \\
К4 & изпълнен & назовани граници (точка~\ref{sec:tm-scope}), демонстриран и затворен пропуск (глава~\ref{sec:case-study}), отбелязан остатъчен риск (точка~\ref{sec:enf-residual}) \\
\bottomrule
\end{tabularx}
\end{table}

Критериите К2 и К4 са изпълнени изцяло. Критериите К1 и К3 са изпълнени за
SQLite, а за сървърните хранилища К1 се подкрепя само от анализа на изходния код.
Разликата между двата вида подкрепа е съществена и не се замъглява: анализът
показва механизма, но само измерването може да покаже, че механизмът работи
така, както е описан.

\subsection{Заплахи за валидността}
\label{sec:eval-validity}

\textbf{Вътрешна валидност.} Оракулът вижда само ефекти и подценява наивния път.
Това изместване работи срещу тезата, а не в нейна полза. Компилационните тестове
използват конектора \code{MockDB}, който превежда заявките без база данни.
Проверката на границата обаче се извършва в общия слой, преди конектора, така че
изборът на конектор не влияе на резултата от тези тестове.

\textbf{Конструктивна валидност.} Диференциалният тест използва една форма на
заявка с две колони. Тя е избрана, за да дава шанс на най-разпространените
полезни товари с \code{UNION}. Друга форма би могла да направи наивния път
уязвим за други полезни товари. Това не засяга резултата на библиотеката,
защото при нея стойността не достига граматиката при никоя форма.

\textbf{Външна валидност.} Измерването е извършено с един диалект и един корпус
от 66 полезни товара. Корпусът не е изчерпателен и съдържа пропуски, които са
изброени в самия него. Твърдението, че гаранцията се отнася за всички входове, не
се основава на корпуса, а на аргумента от точка~\ref{sec:theory-compile-time},
който компилационните тестове проверяват. Корпусът показва, че аргументът издържа
на реални атаки, но не замества аргумента.

\textbf{Надеждност.} Резултатите са детерминирани. Измерванията на времето са
изключени от записаните резултати, за да не предизвикват разлики между
изпълненията. Всеки резултат е записан във файл в изходния код и може да бъде
получен повторно с командите от приложение~А.
